\chapter{System Design}

This chapter describes how Seed Bank is put together: the rules that shape the
backend, the runtime topology, the object model behind the inference pipeline, the
request and authentication flows, the relational schema, the machine-learning
platform that decides which model serves a scan, the synthetic-data tool that
feeds the detector, the client interfaces, and the security posture. The design
follows the principle that the parts of the system that change for different
reasons should be kept apart, so a new inference engine, a new database table, or
a new screen can be added without disturbing the rest \cite{cleanarchitecture2017}.

\section{System Architecture}

\subsection{Five architectural pillars}

The backend is built on five rules that are enforced by code review rather than by
a linter, because they constrain how the layers relate to each other rather than
how any single line is written.

The first rule is that the system is asynchronous from end to end. FastAPI runs on
an asyncio event loop, the database is reached through async SQLAlchemy 2 over
\texttt{asyncpg}, and every external client (Redis, MinIO, ClickHouse, outbound
HTTP) is async \cite{fastapi}. Nothing in the request path blocks the loop; the one
genuinely CPU-bound operation, a PyTorch forward pass, is pushed onto a worker
thread with \texttt{asyncio.to\_thread} so the loop stays responsive.

The second rule is strict layering: \texttt{routers} call \texttt{services}, which
call \texttt{repositories}, which talk to the ORM. Routers parse the request, call
one service, wrap the result in an envelope, and return; they never import
SQLAlchemy. Services hold the business logic and never import FastAPI. Repositories
hold exactly one query or one persistence intent each and never embed business
rules. Each layer can be tested against the one below it without dragging in the
web framework or a live database.

The third rule is that Pydantic validates everything at the boundary. Every request
and response is a typed schema, and the ORM model never leaks past the service
layer, so the shape of an HTTP payload is decoupled from the shape of a database
row. The fourth rule is a single \texttt{Settings} object, built with
pydantic-settings, as the only source of configuration. Values come from
environment variables, a \texttt{.env} file in development, or one file per field
under \texttt{/run/secrets} in production; no code reads \texttt{os.environ}
directly. The fifth rule is model traceability: the chain
\texttt{seed\_detection} $\rightarrow$ \texttt{inference} $\rightarrow$
\texttt{model\_artifact} is a hard foreign-key relationship, so any individual seed
result can be traced back to the exact model version that produced it.

\subsection{Context view}

At the outermost zoom level the platform is a single boundary that three kinds of
actor talk to over HTTPS and JSON. End users (farmers) upload seed images and poll
for results. AI developers register and promote models and run offline
experiments. Administrators manage users, suppliers, and traffic splits. All three
hit the same HTTP surface; what differs is the role each one carries and the
operations that role is allowed to perform. Outside the boundary sit two OAuth
identity providers (Google and GitHub), an optional hosted Roboflow inference
backend, and an optional Sentry endpoint for error tracking. Model training
happens outside the system entirely: only the trained weights and their metrics
are imported, through \texttt{scripts/register\_model.py} or \texttt{POST /models}.
\autoref{fig:context} shows this view.

\figslot{01-system-context.pdf}{System context: the actors and external systems that interact with the Seed Bank boundary.}{fig:context}

\subsection{Container view}

One level deeper, the platform is a set of containers declared in a single
\texttt{compose.yaml}, each one a process with a healthcheck and explicit CPU and
memory limits \cite{merkel2014docker}. The \texttt{api} container serves every
router under \texttt{/api/v1}. Two Celery worker containers consume disjoint sets
of queues: \texttt{worker-inference} runs the detect-then-classify pipeline and
offline evaluation, while \texttt{worker-cpu} handles data-warehouse writes,
experiment bookkeeping, and housekeeping. The split exists so the heavy ML
dependencies (\texttt{torch}, \texttt{torchvision}, around 1.6\,GB) load only into
the inference worker and never weigh down a routine warehouse write. Four stateful
services back the application: PostgreSQL for the OLTP database, Redis for cache,
Celery broker, and rate-limit store, MinIO for object storage (images, weights,
datasets, experiment reports), and ClickHouse for the analytics warehouse. MLflow
tracks experiments. Every host port is bound to \texttt{127.0.0.1}, so nothing is
exposed to the local network by default. \autoref{fig:containers} shows the stack.

\figslot{02-containers.pdf}{Container view of the Compose stack: application processes, stateful services, and ML tooling.}{fig:containers}

\subsection{Component view}

Inside the \texttt{api} container the layering becomes visible as concrete classes.
A request enters through middleware (request-id assignment, structured logging,
CORS, the rate limiter, and the problem-details exception handler), reaches a
versioned router under \texttt{api/v1/}, and is wired to its service through
\texttt{api/deps.py}, which resolves the current user, checks the required role,
and supplies the database session and repositories. Each router maps to one
service: \texttt{auth.py} to \texttt{AuthService}, \texttt{analyze.py} to
\texttt{AnalysisService}, \texttt{batches.py} to \texttt{BatchService},
\texttt{models.py} and \texttt{traffic.py} to \texttt{ModelRegistryService} and
\texttt{TrafficRouter}, and so on. Services reach the database only through
repositories such as \texttt{UserRepository}, \texttt{ScanBatchRepository}, and
\texttt{SeedDetectionRepository}, and reach the outside world through
infrastructure adapters for storage, cache, OAuth, ML, and analytics. The ORM, a
single SQLAlchemy 2 \texttt{AsyncSession}, sits at the bottom \cite{sqlalchemyguide}.
\autoref{fig:api-components} shows the layering, and Table~\ref{tab:ch4-errmap}
shows how each domain exception that a service may raise is translated into an HTTP
status code.

\figslot{03-api-components.pdf}{API component layering inside the \texttt{api} container: middleware, routers, services, repositories, and adapters.}{fig:api-components}

\begin{singlespace}
\begin{table}[H]
\centering
\caption{Mapping from domain exception to HTTP status and problem-details code.}
\label{tab:ch4-errmap}
\begin{tabular}{lll}
\toprule
Domain exception & HTTP status & \texttt{code} field \\
\midrule
\texttt{ValidationError} & 422 & \texttt{validation\_error} \\
\texttt{AuthError} & 401 & \texttt{auth\_error} \\
\texttt{ForbiddenError} & 403 & \texttt{forbidden} \\
\texttt{NotFoundError} & 404 & \texttt{not\_found} \\
\texttt{ConflictError} & 409 & \texttt{conflict} \\
\texttt{RateLimitError} & 429 & \texttt{rate\_limited} \\
\texttt{ExternalServiceError} & 502 / 503 & \texttt{external\_service\_error} \\
\bottomrule
\end{tabular}
\end{table}
\end{singlespace}

The worker container has its own component structure. The Celery app factory routes
the \texttt{seedbank.analyze\_image} task to the \texttt{inference} queue with
\texttt{acks\_late} and a prefetch of one, so a worker only acknowledges a task
after it succeeds or finally fails, and processes one image at a time. The task is
a synchronous entry point that calls \texttt{asyncio.run} on the real async work.
Each task opens a fresh \texttt{AsyncEngine} and session through
\texttt{worker\_session\_scope} and disposes it at the end. This is not an
optimisation oversight: \texttt{asyncpg}'s connection pool binds to the event loop
it was created on, and every \texttt{asyncio.run} creates a new loop, so reusing
the API's process-wide engine would fail with a future attached to a different
loop. Within one task the worker runs a fixed sequence: compare-and-set the batch
from pending to running, fetch the image bytes from MinIO, resolve the detection
model through the traffic router, detect, persist the inference and its detections,
commit, resolve a classifier, crop and classify each detection, update the quality
column, commit again, then finalise the batch. \autoref{fig:worker-components}
shows these components.

\figslot{04-worker-components.pdf}{Inference worker components: the Celery task, its per-task session scope, and the detect-then-classify steps.}{fig:worker-components}

\section{Class Diagrams}

The object model behind the ML pipeline is designed so that the rest of the system
never imports PyTorch. The seam that makes this possible is a duck-typed
\texttt{Protocol}, \texttt{InferenceBackend}, with two methods:
\texttt{detect(image, DetectionConfig)} returns a list of \texttt{Detection}, and
\texttt{classify(crop, ClassificationConfig)} returns a \texttt{Classification}.
Both return framework-free data-transfer objects. A \texttt{Detection} carries a
\texttt{BoundingBox}, a confidence as a \texttt{Decimal}, and a class name; a
\texttt{Classification} carries a label and a confidence. Because these DTOs hold
no tensors, torch never leaks up through the service and router layers, which is
why the API image can be built without it.

Three classes satisfy the protocol. \texttt{TorchLocalBackend}, the default, loads
first-party \texttt{.pth} weights from MinIO and runs a local PyTorch forward pass.
\texttt{YoloBackend} wraps the Ultralytics package for the YOLO family
\cite{yolov8}, and \texttt{RoboflowBackend} calls a hosted Roboflow endpoint over
HTTP. Adding a fourth backend means writing one class that matches the protocol and
registering it in the pipeline factory; no caller changes. This is the strategy
pattern applied to the inference engine \cite{designpatterns1994}.

Model architectures are kept separate from the engines that run them. A builder is a
factory function registered under a \texttt{builder\_key} that constructs the bare
\texttt{nn.Module}; weights are loaded into it afterwards. The
\texttt{ModelBuilder} registry maps a key to its builder, so a new architecture is
one new file with a decorator and changes neither the backends nor the router.
Builders are append-only by convention: when the math of an architecture changes,
the team copies it to a new \texttt{-vN} key rather than mutating a builder that a
production model already references, because mutating it would silently change what
a deployed model computes.

The expensive part, the loaded weights, lives in \texttt{ModelManager}, one
instance per worker process. It holds a process-wide LRU cache of loaded modules
(default four, with separate caches for torch and YOLO). When a model is requested
it builds the module through the registry, fetches the weights from MinIO, loads the
state dict, moves it to GPU or CPU, and switches it to evaluation mode. Concurrent
loads of the same \texttt{model\_id} are serialised through a per-id
\texttt{asyncio.Lock} so two tasks never duplicate one GPU load, the least-recently
used model is evicted to bound memory, and a model is hot-reloaded when its
\texttt{model\_artifacts.updated\_at} advances.

These pieces are coordinated by two services. \texttt{TrafficRouter} answers the
question of which model should run for a given segment, and \texttt{AnalysisService}
drives the request: it persists the batch and its images through
\texttt{ScanBatchRepository} and \texttt{ScanImageRepository}, which expose intents
like \texttt{cas\_status} for the concurrency-safe state transitions. The same
service-then-repository layering described for the API holds throughout: a service
never issues raw SQL, and a repository never decides policy. \autoref{fig:class}
shows the key classes and their relationships.

\figslot{13-class.pdf}{Key classes in the ML pipeline: the \texttt{InferenceBackend} protocol, its three backends, the builder registry, the model manager, and the services that use them.}{fig:class}

\section{Sequence Diagrams}

\subsection{The analyze flow}

The analyze flow is the path the rest of the platform exists to serve. A client
sends a multipart \texttt{POST /analyze} carrying one or more image files and
optional metadata (seed type, supplier, country, GPS, and, for developers, a model
override). The ordering inside \texttt{AnalysisService} is load-bearing. First it
validates every file: the count must be at most sixteen, each must pass size and
MIME checks, and each must decode as a real image through Pillow. Then it writes
the image bytes to MinIO \emph{before} the database commit, so any committed row
always references an object that actually exists. Next it creates the
\texttt{scan\_batch} in the \texttt{pending} state along with one
\texttt{scan\_image} per file, writes an audit-log entry, and commits. Only after
the commit does it dispatch one Celery task per image onto the inference queue, so a
worker can never pick up a task for a batch that is not yet visible in the database.
The endpoint returns \texttt{202 Accepted} with a \texttt{Location} header pointing
at the new batch.

Each worker task then runs the per-image pipeline. It compare-and-sets the batch
from \texttt{pending} to \texttt{running}, which only the first task to arrive wins,
so several workers processing a multi-image batch cannot corrupt the state. It
fetches the image from MinIO, resolves the detection model, runs detection in a
thread, and persists one \texttt{inference} row and N \texttt{seed\_detections} rows
with normalized bounding boxes. After committing the detection results it resolves a
classifier, groups the detections by seed type, crops each detected seed from the
source image, classifies it as good or bad, and bulk-updates the quality column. If
no classifier is registered for a detected type, those detections are simply left
unclassified rather than treated as an error. Finally, when every image in the batch
has a detection inference, it compare-and-sets the batch to \texttt{succeeded},
\texttt{partial}, or \texttt{failed}. The client meanwhile polls
\texttt{GET /batches/\{id\}} roughly every two seconds until the batch reaches a
terminal status, at which point a single eager-loaded query returns the batch with
its images, inferences, and detections. \autoref{fig:analyze-seq} shows the
submission, worker, and polling phases.

\figslot{06-analyze-sequence.pdf}{Analyze request sequence: submission and dispatch, per-image worker consumption, and client polling.}{fig:analyze-seq}

Crash safety follows from this ordering. A transient MinIO or Redis fault raises
\texttt{ExternalServiceError}, which Celery auto-retries up to twice; validation and
not-found errors are never retried. Because detection rows are committed before
classification runs, a failure during classification degrades the batch to
\texttt{partial} and keeps the detection data, rather than throwing away good
results. If no model can be routed at all, the batch fails with a human-readable
message that the clients display, so a user with an unprovisioned deployment sees a
clear explanation rather than a blank failure.

\subsection{Authentication flows}

Three authentication flows live behind \texttt{/api/v1/auth/}. In local login the
service looks up the user by email, verifies the password against the bcrypt hash,
and on success inserts a refresh-token row (storing only a SHA-256 hash of the
token) and signs a short-lived access JWT carrying the role and scopes
\cite{rfc7519jwt}. When the access token later expires the client calls
\texttt{POST /auth/refresh}. The service looks up the refresh token by its hash; if
the token is revoked or expired it rejects the request, and if it is valid it
revokes the old row, records the new token as its replacement in a rotation chain,
and issues a fresh pair. Because each refresh rotates the token and links old to
new, presenting a previously used refresh token is detectable as replay and
invalidates the chain.

The OAuth flow uses the standard authorization-code exchange
\cite{rfc6749oauth}. The API redirects the browser to the provider's authorize URL
with a \texttt{state} value carried across the redirect by the session middleware,
receives the callback, exchanges the code for tokens, fetches the user's email and
subject, and either finds the linked account or creates a user and links it. It
then issues the same token pair as local login. A third path covers headless
clients: an API key presented as a bearer credential is recognised by its prefix,
short-circuits JWT verification, is hashed and looked up, and yields an
authenticated identity with the role and scopes attached to the key.
\autoref{fig:auth-seq} shows the login-and-rotation, OAuth, and API-key sequences.

\figslot[0.85]{08-auth-sequence.pdf}{Authentication sequences: local login with refresh-token rotation, OAuth code flow, and per-request API-key authentication.}{fig:auth-seq}

\section{Database Entity-Relationship Diagram}

The OLTP store is a PostgreSQL schema of eighteen tables, managed with Alembic
\cite{postgresinprod}. Every primary key is a UUIDv7 generated in application code,
which keeps keys time-ordered and therefore index-friendly while staying globally
unique. The tables fall into four groups.

The identity and authentication group has five tables: \texttt{users} (email,
bcrypt hash, role, verification and active flags), \texttt{refresh\_tokens} (token
hash, expiry, revocation, and the \texttt{replaced\_by\_id} that forms the rotation
chain), \texttt{oauth\_accounts} (provider and subject, encrypted provider tokens),
\texttt{api\_keys} (key hash, prefix, scopes, expiry), and an append-only
\texttt{audit\_log} indexed by actor and time. The reference group has two tables,
\texttt{seed\_types} and \texttt{suppliers}, that supply catalog data.

The inference graph is the heart of the schema and the four tables that carry the
traceability chain. A \texttt{scan\_batches} row is created per
\texttt{POST /analyze} and holds the state machine
(\texttt{pending} $\rightarrow$ \texttt{running} $\rightarrow$
\texttt{succeeded} / \texttt{partial} / \texttt{failed}), timing, and optional geo
metadata. Each batch groups several \texttt{scan\_images}, each of which records its
MinIO storage key, content hash, and pixel dimensions. An \texttt{inferences} row
is one model run over one image and records the backend, the \texttt{model\_id}, the
latency, and any error. Each detected seed becomes a \texttt{seed\_detections} row
with its normalized bounding box, confidence, and quality, linked back to its
inference. The data fans out as one image to N detections to N quality labels.

The ML-platform group has seven tables: \texttt{model\_artifacts} (registered
weights plus metadata and lifecycle status), \texttt{model\_metrics} (per-model
performance records), \texttt{traffic\_splits} (weighted A/B routing with a temporal
validity window), \texttt{datasets} and \texttt{dataset\_items} (offline-evaluation
ground truth), and \texttt{experiments} and \texttt{experiment\_results}
(offline-evaluation runs and their computed metrics). \autoref{fig:erd} shows the
full schema, and Table~\ref{tab:ch4-invariants} lists the invariants the design
relies on.

\figslot[0.85]{05-db-erd.pdf}{Database entity-relationship diagram: identity, reference, the inference traceability graph, and the ML-platform tables.}{fig:erd}

\begin{singlespace}
\begin{table}[H]
\centering
\caption{Schema invariants enforced by the database design.}
\label{tab:ch4-invariants}
\begin{tabular}{p{0.34\textwidth}p{0.58\textwidth}}
\toprule
Invariant & What it guarantees \\
\midrule
Traceability chain (FK) & Every \texttt{seed\_detections} row points to one \texttt{inferences} row, which points to exactly one \texttt{model\_artifacts} row. \\
\texttt{uq\_inferences\_image\_model} & One inference row per \texttt{(image\_id, model\_id)}; a re-run with the same model updates rather than duplicates. \\
\texttt{uq\_scan\_images\_batch\_storage\_key} & Re-uploading the same bytes in the same batch fails fast, because the storage key embeds the SHA-256. \\
Normalized bounding boxes & Box columns are \texttt{NUMERIC(7,6)} in the range $[0,1]$; pixel coordinates are derived at render time, never stored. \\
Decimal confidences & \texttt{confidence} is \texttt{NUMERIC(5,4)}, so arithmetic over the audit trail does not lose precision to float drift. \\
Append-only audit log & \texttt{audit\_log} is never updated or soft-deleted; the temporal validity window on \texttt{traffic\_splits} makes promotions inserts, not destructive updates. \\
\bottomrule
\end{tabular}
\end{table}
\end{singlespace}

Storing bounding boxes normalized to $[0,1]$ as the source of truth has a direct
runtime consequence: crops are derived, never stored. The worker multiplies a
normalized box by the source image's pixel dimensions at classify time to cut each
crop, so the same detection survives image resizing and the database never holds an
image. Confidences and box coordinates are also emitted to clients as decimal
strings rather than floats, and parsed to numbers only at render time, for the same
precision reason.

\section{ML Platform Design}

The ML platform turns a trained \texttt{.pth} file into production traffic and lets
developers experiment against production without disturbing it. It is organised
around three ideas: a registry with a lifecycle, a router that selects a model per
request, and a manager that loads and caches weights.

\subsection{Registry and lifecycle}

Weights are uploaded to MinIO and recorded in \texttt{model\_artifacts}, either
through \texttt{scripts/register\_model.py} or \texttt{POST /models}. Each artifact
carries its kind (detection or classification), its backend, the MinIO key, a
\texttt{builder\_key}, a free-form \texttt{config} JSON of per-model knobs
(confidence and IoU thresholds, maximum detections, image size, classification
threshold), an optional seed type, and a lifecycle status. The status moves through
\texttt{registered} $\rightarrow$ \texttt{staging} $\rightarrow$ \texttt{production}
$\rightarrow$ \texttt{archived} by an explicit API action, and a model can also be
archived directly from \texttt{registered} or \texttt{staging} when a promotion is
abandoned. Only \texttt{staging} and \texttt{production} models are eligible to
serve requests, which is the same condition a developer override must satisfy. This
registry plays the role that MLflow plays for tracking, but for serving: MLflow
records the offline experiment runs, while the registry governs what actually runs
in production \cite{mlflow2018}.

\subsection{Traffic router}

Given a segment \texttt{(kind, seed\_type\_id)}, the router decides which model
runs. It first reads the active \texttt{traffic\_splits} rows for the segment, cached
in Redis for sixty seconds \cite{redisinaction}. If splits exist it routes by a
sticky bucket, \texttt{hash(user\_id) \% 100}, against the cumulative weights, so the
same user always lands on the same model. That stickiness is what makes an A/B test
stable: a user does not flip between models on consecutive scans, yet traffic is
spread uniformly across the user base. If no splits are configured the router falls
back to the production model for that seed type, and then to the global,
seed-type-agnostic production model. This fallback makes per-type promotion
optional: a deployment can promote one global model and every scan routes to it,
including the mobile point-and-shoot flow, which sends no seed type at all. If
nothing is routable the router raises \texttt{ModelNotReadyError}, which becomes the
clear failure message the clients show.

A developer or administrator can override the router with an explicit
\texttt{model\_id} at analyze time. The override must name a \texttt{staging} or
\texttt{production} model of the requested kind, and an end user attempting it is
rejected with \texttt{ForbiddenError}. This override is the platform's escape hatch
for trying one model on one image without touching production traffic, and the
inference row still records the overridden \texttt{model\_id}, so the result remains
traceable. \autoref{fig:ml-platform} shows the lifecycle, the router's decision
tree, and the override path.

\figslot{09-ml-platform.pdf}{Model lifecycle, weighted traffic splits, and the per-request developer override.}{fig:ml-platform}

\subsection{Model manager}

The router decides which model; the manager loads it. Each worker process holds one
\texttt{ModelManager} with a process-wide LRU cache, a per-id load lock, and a
hot-reload check against \texttt{updated\_at}, as described in the class-diagram
section. The two wired architectures show why the builder and backend seams matter.
Detection uses a Faster R-CNN with a ResNet-50 backbone and a Feature Pyramid
Network and a three-class head over background, coffee, and maize
\cite{fasterrcnn2015,resnet2015,lin2017fpn}; the backend filters by the configured
confidence threshold, normalizes the pixel boxes, and caps the count.
Classification uses a ResNet-18 backbone with a Convolutional Block Attention Module
inserted after the last residual stage, then hybrid average-and-max pooling into a
single-logit head \cite{cbam2018}. There is one classifier per crop type, which is
why detections are grouped by seed type before classification. Because a detector
and a classifier are recorded as separate inference rows over the same image, each
can be versioned, swapped, and A/B-tested independently.

\section{MultiSeedGen Design}

Training a seed detector needs many photographs of cluttered, overlapping seeds
with a correct bounding box on every seed, which are slow to capture and tedious to
annotate by hand. MultiSeedGen is the companion tool that manufactures this data: it
takes photographs of single seeds and composites them into annotated multi-seed
scenes whose labels are exact by construction, an instance of cut-and-paste
synthesis \cite{dwibedi2017cutpaste}. The whole pipeline is driven by one flat
Pydantic configuration object, and its output is byte-identical for a fixed
configuration and seed, which lets a generated dataset be reproduced exactly.

The pipeline runs in stages, shown in \autoref{fig:msg-pipeline}. It begins with a
segmentation cascade that cuts each seed out of its source photo: a classical
colour-and-brightness method (border colour, then Otsu thresholding, then GrabCut)
is tried first, and a learned matting model is the fallback, either U2-Net through
rembg \cite{qin2020u2net} or Segment Anything \cite{kirillov2023sam}. A confidence
gate decides whether the cut-out is clean enough, and every result is stored in a
content-hashed disk cache so a seed is segmented once and reused across runs. The
clean cut-outs form a seed pool.

Scene composition then paints seeds onto a canvas. It samples a number of seeds per
scene, places each one while rejecting any placement whose box overlaps an existing
box beyond a configured IoU limit and that fails a visibility check, applies
label-safe geometric and photometric augmentation (scale, rotation, shear, and HSV
jitter chosen so they do not move a label box), composites onto a background that is
solid, gradient, noise, texture, or a real inpainted photo, and adds soft
directional shadows. Because the tool knows exactly where it placed each seed, label
derivation is occlusion-aware: it produces axis-aligned bounding boxes and polygons
that account for seeds partly hidden behind others, and records provenance for each
box back to its source seed. The composed scene then passes through camera-matched
degradation (blur, noise, JPEG compression, and gamma or vignette) so the synthetic
images resemble real phone photos, a form of domain randomization that helps a model
trained on them transfer to real input \cite{tobin2017domainrand,shorten2019augmentation}.

Export writes the dataset in the formats a detector trainer expects: YOLO boxes,
YOLO-seg polygons, COCO annotations \cite{lin2014coco}, and a \texttt{data.yaml} plus
a run manifest. The final stage is a leakage-safe split: a source seed reserved for
validation never appears in any training image, enforced per class and
deterministically, so the validation score is not inflated by the same seed
appearing on both sides of the split. The resulting dataset is what trains the Seed
Bank detection model.

\figslot{14-multiseedgen-pipeline.pdf}{MultiSeedGen synthetic-data pipeline: segmentation cascade, scene composition, occlusion-aware labels, degradation, export, and leakage-safe split.}{fig:msg-pipeline}

\section{User Interface Design (Wireframes and Mockups)}

Seed Bank ships two clients that share a contract and a visual language. The web app
is a React 18 single-page app served by Vite, and the mobile app is built with Expo
and React Native \cite{reactindepth,reactnativeexpo}. Both speak the same
\texttt{/api/v1} paths, the same \texttt{\{data\}} response envelopes, and the same
bearer authentication with a single transparent refresh on a 401. Both wear an
agricultural-green theme, a deep leaf-green primary with warm amber accents over a
warm off-white field in light mode and a soil-night palette in dark mode, defined as
HSL CSS variables so light and dark swap without a rebuild. Both are bilingual:
English and Arabic, with Arabic flipping the entire layout to right-to-left.

\subsection{Web app}

The web app opens on a sign-in screen with a field-pattern backdrop and theme and
language toggles available before authentication, shown in \autoref{fig:web-login}.
After sign-in the dashboard (\autoref{fig:web-dashboard}) presents a strip of
key figures (scans, images, success rate, and a fourteen-day sparkline), quick
actions, and recent scans. The core journey runs through the analyze page
(\autoref{fig:web-analyze}), where a user drags in or picks photos, optionally adds
metadata, and submits, which starts a batch and navigates to its detail. The batch
detail page (\autoref{fig:web-batch-detail}) polls until the batch is terminal, then
shows an insights panel (a quality donut, a confidence histogram, and headline
stats), a per-seed-type breakdown, and per-image cards with an interactive
bounding-box overlay that can filter by quality, threshold by confidence, and toggle
labels. From here a user can export the results as CSV or JSON, or create a public
read-only share link.

\figslot{web-login.png}{Web sign-in, with theme and language toggles available before authentication.}{fig:web-login}

\figslot{web-dashboard.png}{Web dashboard, showing the key-figure strip, quick actions, and recent scans.}{fig:web-dashboard}

\figslot{web-analyze.png}{Web analyze page, where photos and optional metadata are submitted to start a batch.}{fig:web-analyze}

\figslot{web-batch-detail.png}{Web batch detail with the interactive bounding-box overlay and the insights panel.}{fig:web-batch-detail}

Right-to-left support is achieved by converting directional Tailwind utilities to
logical properties in the shared primitives and pages, with direction-aware chevrons
and a drawer that opens from the correct edge. Numbers stay in Latin digits in
Arabic for cross-script consistency, and dates are formatted with the active locale.
The entire end-user surface is translated; the role-gated admin and ML pages are
deliberately left in English, since farmers never reach them.

\subsection{Mobile app}

The mobile app is built for farmers in the field, and its headline screen is the
camera (\autoref{fig:mobile-camera}): a live preview with a centre framing guide,
a shutter that captures multiple shots into a review strip, and a ``Check seeds''
button that uploads the captured photos as one batch. The result screen
(\autoref{fig:mobile-result}) then polls the batch until it is terminal and shows
the good-rate, the seed counts, and the confidence. The same capture-and-submit
design extends to a production line: up to sixteen frames of seeds passing under a
fixed camera accumulate in the review strip and are submitted as one batch, then
graded in parallel by the inference workers, so an operator gets a report on the
whole pass within seconds without any belt-specific code.

\figslot[0.45]{mobile-camera.png}{Mobile camera capture, with a live preview, framing guide, and multi-shot review strip.}{fig:mobile-camera}

\figslot[0.45]{mobile-result.png}{Mobile result screen, polling the batch and showing good-rate, counts, and confidence.}{fig:mobile-result}

\subsection{MultiSeedGen interface}

MultiSeedGen has its own web interface, shown in \autoref{fig:msg-ui}: a React
single-page app served by a thin FastAPI wrapper that drives the same command-line
pipeline as a subprocess and streams its log over a websocket. The interface is
organised into four tabs, for running a generation, tuning the segmentation, browsing
generated datasets, and editing the configuration. The configuration form is built
from the same Pydantic model that defines the CLI, so the form fields, their choices,
and their defaults can never drift from what the pipeline actually accepts.

\figslot{msg-ui.png}{MultiSeedGen web UI, with tabs for running, segmentation tuning, dataset browsing, and configuration.}{fig:msg-ui}

\section{Security Design Considerations}

Security is built into the layers rather than bolted on at the edge, and the design
choices below cover authentication, authorization, transport, and error handling.

Local authentication stores passwords only as bcrypt hashes; a plaintext password is
never written anywhere. Sessions use two JSON Web Tokens, a short-lived access token
(fifteen minutes) and a long-lived refresh token (days), and the refresh token
rotates on every use \cite{rfc7519jwt}. Only the SHA-256 hash of a refresh token is
stored, each rotation links the old token to its replacement, and presenting a
previously used token is treated as replay and invalidates the whole chain. Social
login is offered through Google and GitHub over the OAuth authorization-code flow,
with the \texttt{state} parameter carried across the redirect to guard against
cross-site request forgery on the callback \cite{rfc6749oauth}. API keys give
headless clients a credential that is hashed at rest, carries a recognisable prefix,
and is shown in plaintext exactly once at creation; like refresh tokens, the stored
form is a hash, so a database leak does not expose a usable key.

Authorization is role-based across three roles. An \texttt{end\_user} can analyze
and see only their own history; an \texttt{ai\_developer} adds model, dataset, and
experiment access and the per-request model override; and an \texttt{admin} adds
traffic splits and user management. Roles are checked at the API through a
\texttt{require\_role} dependency and again in the clients' routing, and the model
override is the most heavily tested authorization boundary in the system because it
is the one place a non-admin role gains a privileged capability. Each route also
carries its own rate limit, backed by Redis through slowapi, with tighter caps on
login, registration, refresh, and analyze than on the general default.

Every sensitive action is written to the append-only audit log, which records the
actor, the action, the target, and the time, and is never updated or deleted.
Secrets are read from environment variables in development but from one file per
field under \texttt{/run/secrets} in production, and a pre-commit secret scanner
keeps credentials out of version control. Errors are returned as RFC 9457
problem-details JSON with a correlating request id and never as HTML pages, so an
error surface leaks no stack traces \cite{rfc9457problem}. The deployment binds every
host port to the loopback interface, restricts cross-origin requests to a configured
allow-list, and enforces a trusted-hosts list, so the attack surface exposed by
default is small.
